\chapter{Implementation Details}
\label{chap:implementation}

This chapter describes the software implementation, directory structure, and development workflow.

\section{Software Architecture}

\subsection{Directory Structure}

\begin{lstlisting}[language=bash]
noma_gnn_backup/
├── config.py                  # Global configuration
├── requirements.txt           # Python dependencies
├── data/
│   ├── dataset.py            # PyTorch Geometric dataset
│   ├── normalization.py      # Feature scaling
│   ├── raw/                  # CSV channel data
│   └── processed/            # Processed graphs
├── models/
│   └── pairpower_gnn.py      # NOMA-GNN architecture
├── training/
│   ├── train.py              # Training loop
│   └── losses.py             # Multi-task loss
├── inference/
│   └── infer_pairing.py      # Real-time inference
├── scripts/
│   ├── prepare_data.py       # Dataset generation
│   └── verify_sic_pairs.py   # SIC validation
├── utils/
│   ├── matching.py           # Greedy matching
│   └── metrics.py            # Evaluation metrics
└── checkpoints/
    └── best_model.pt         # Trained weights
\end{lstlisting}

\subsection{Technology Stack}

\begin{table}[H]
\centering
\begin{tabular}{ll}
\toprule
\textbf{Component} & \textbf{Technology} \\
\midrule
Language & Python 3.9 \\
Deep Learning & PyTorch 1.13 \\
Graph Processing & PyTorch Geometric 2.2 \\
Numerical Computing & NumPy 1.23, SciPy 1.9 \\
Graph Algorithms & NetworkX 2.8 \\
Visualization & Matplotlib 3.6, Seaborn 0.12 \\
Progress Tracking & tqdm 4.64 \\
Hardware & NVIDIA RTX 3060 (6GB) \\
OS & Ubuntu 20.04 / Windows 10 \\
\bottomrule
\end{tabular}
\caption{Software Stack}
\end{table}

\section{Key Modules}

\subsection{Dataset Generation (scripts/prepare_data.py)}

Generates 3GPP-compliant channel realizations and constructs graph datasets.

\textbf{Key Functions:}
\begin{itemize}
    \item \texttt{generate_uma_channels()}: Implements 3GPP UMa model
    \item \texttt{construct_candidate_graph()}: Builds edge list with constraints
    \item \texttt{compute_ground_truth()}: Runs Blossom for labels
    \item \texttt{save_to_csv()}: Exports dataset
\end{itemize}

\subsection{Graph Dataset (data/dataset.py)}

PyTorch Geometric dataset loader.

\begin{lstlisting}[language=Python, caption={Dataset Class}]
from torch_geometric.data import Dataset, Data

class NOMAGraphDataset(Dataset):
    def __init__(self, root, transform=None):
        super().__init__(root, transform)
        self.data_list = self.process()
    
    def len(self):
        return len(self.data_list)
    
    def get(self, idx):
        return self.data_list[idx]
    
    def process(self):
        # Load CSV files
        # Convert to PyG Data objects
        # Return list of graphs
        pass
\end{lstlisting}

\subsection{Model (models/pairpower_gnn.py)}

Complete NOMA-GNN implementation.

\begin{lstlisting}[language=Python, caption={Model Class Skeleton}]
import torch
import torch.nn as nn
from torch_geometric.nn import SAGEConv

class PairPowerGNN(nn.Module):
    def __init__(self, in_dim=5, hidden_dim=128, num_layers=3):
        super().__init__()
        # Encoder
        self.input_proj = nn.Linear(in_dim, hidden_dim)
        self.convs = nn.ModuleList([
            SAGEConv(hidden_dim, hidden_dim) 
            for _ in range(num_layers)
        ])
        
        # Decoders
        self.pair_mlp = self._build_mlp(256, [128, 64, 1])
        self.rate_mlp = self._build_mlp(256, [128, 64, 1])
        self.alpha_mlp = self._build_mlp(256, [128, 64, 1])
    
    def forward(self, data):
        # Encode nodes
        x = F.relu(self.input_proj(data.x))
        for conv in self.convs:
            x = F.relu(conv(x, data.edge_index))
        
        # Decode edges
        edge_emb = torch.cat([x[data.edge_index[0]], 
                             x[data.edge_index[1]]], dim=1)
        
        p_pair = torch.sigmoid(self.pair_mlp(edge_emb))
        R_sum = F.relu(self.rate_mlp(edge_emb))
        alpha = torch.sigmoid(self.alpha_mlp(edge_emb))
        
        return p_pair, R_sum, alpha
\end{lstlisting}

\subsection{Training (training/train.py)}

Main training script with checkpointing and logging.

\subsection{Inference (inference/infer_pairing.py)}

Real-time pairing for deployed systems.

\section{Configuration Management}

Global settings in \texttt{config.py}:

\begin{lstlisting}[language=Python, caption={Configuration}]
from dataclasses import dataclass

@dataclass
class Config:
    # System parameters
    TOTAL_POWER: float = 1.0
    NOISE_POWER: float = 1e-9
    SIC_THRESHOLD_DB: float = 8.0
    BANDWIDTH_HZ: float = 20e6
    
    # Graph construction
    USE_ANGLE_GUARD: bool = True
    MIN_ANGLE_DEG: float = 25.0
    
    # Model architecture
    IN_CHANNELS: int = 5
    HIDDEN_DIM: int = 128
    NUM_LAYERS: int = 3
    DROPOUT: float = 0.2
    
    # Training
    LR: float = 1e-3
    WEIGHT_DECAY: float = 1e-4
    EPOCHS: int = 200
    BATCH_SIZE: int = 32
    
    # Loss weights
    LAMBDA_BCE: float = 1.0
    LAMBDA_RSUM: float = 0.5
    LAMBDA_ALPHA: float = 0.3

CFG = Config()
\end{lstlisting}

\section{Development Workflow}

\subsection{Data Preparation}
\begin{lstlisting}[language=bash]
# Generate 50k scenarios
python scripts/prepare_data.py --num_scenarios 50000 --output data/raw/
\end{lstlisting}

\subsection{Training}
\begin{lstlisting}[language=bash]
# Train model
python training/train.py --config config.py --gpu 0
\end{lstlisting}

\subsection{Evaluation}
\begin{lstlisting}[language=bash]
# Test on new scenarios
python inference/infer_pairing.py --checkpoint checkpoints/best_model.pt
\end{lstlisting}

\section{Code Quality}

\subsection{Testing}
Unit tests for critical components:
\begin{itemize}
    \item SIC constraint verification
    \item Graph construction validity
    \item Power allocation correctness
    \item Rate calculation accuracy
\end{itemize}

\subsection{Documentation}
\begin{itemize}
    \item Inline docstrings (NumPy style)
    \item Type hints for function signatures
    \item README files in each subdirectory
\end{itemize}

\section{Reproducibility}

All experiments are reproducible via:
\begin{itemize}
    \item Fixed random seeds (42)
    \item Version-controlled dependencies (requirements.txt)
    \item Saved hyperparameters (config.py)
    \item Checkpointed models (best\_model.pt)
\end{itemize}

\section{Performance Optimization}

\subsection{GPU Acceleration}
\begin{itemize}
    \item All tensor operations on CUDA
    \item Mini-batch processing
    \item Mixed-precision training (optional)
\end{itemize}

\subsection{Memory Efficiency}
\begin{itemize}
    \item Lazy dataset loading
    \item Gradient checkpointing for large graphs
    \item Sparse edge storage
\end{itemize}
